\subsection{分辨率折叠映射（folding）}\label{subsec:folding-map}

折叠的目标在于将任意窗口微态 $\omega=\omega_1\cdots \omega_m\in\Omega_m$ 规范化为一个合法词 $x\in X_m$。

为避免“人为编号”，本文直接以 Fibonacci 权重赋予微态可审计整数值，再取 Zeckendorf 规范形：
\begin{definition}[Fibonacci 加权值与折叠算子]\label{def:fold-word}
对 $\omega=\omega_1\cdots\omega_m\in\Omega_m$ 定义其 Fibonacci 加权值
\[
N(\omega):=\sum_{k=1}^{m}\omega_k F_{k+1}\ \in\ \NN.
\]
令 $Z(N(\omega))=(c_1,c_2,\dots)$ 为 $N(\omega)$ 的 Zeckendorf 数字（\cite{Zeckendorf1972}），定义分辨率 $m$ 的折叠
\[
\Fold_m:\Omega_m\to X_m,\qquad
\Fold_m(\omega):=\pi_m\!\bigl(Z(N(\omega))\bigr).
\]
\leanverified{Fold}
\end{definition}

\begin{remark}[折叠映射应读作“合法构造”，不只是取商]\label{rem:fold-map-admissible-construction}
定义 \ref{def:fold-word} 的重点不只是把 \(\Omega_m\) 压到 \(X_m\)，而是给出一条可审计的对象生成链：先取 Fibonacci 加权值，再进入 Zeckendorf 规范形，再截断到分辨率 \(m\) 的可见层。按 Euclid 式口径，这一步属于构造层，而不是刚性层。后面的幂等、满射、纤维规模与残差定理，才是在“对象已被合法构造出来”之后，进一步描述它有多唯一、还剩多少不可见自由度。
\end{remark}

\subsubsection{截断 Zeckendorf 位的模分解与顶位统计}\label{subsubsec:fold-zeckendorf-mod-topbits}

在定义 \ref{def:fold-word} 的记号下，记
\[
v_m(x):=\sum_{k=1}^{m}x_kF_{k+1}\qquad(x\in X_m),
\qquad
M_m:=F_{m+2},
\qquad
T_m:=\sum_{k=1}^{m}F_{k+1}=F_{m+3}-2.
\]
并令 $Z(N(\omega))=(c_1,c_2,\dots)\in X_\infty$ 为 Zeckendorf 数字序列；为强调“顶位”含义，下文记
\[
(\zeta_m,\zeta_{m+1}):=(c_m,c_{m+1})\in\{(0,0),(1,0),(0,1)\}.
\]

\begin{theorem}[截断位的模分解与唯一性]\label{thm:fold-zeckendorf-mod-decomposition-unique}
对任意 $m\ge 1$ 与任意 $\omega\in\Omega_m$，存在且仅存在一对 $(x,c)\in X_m\times\{0,1\}$ 使
\[
N(\omega)=v_m(x)+c\,M_m,
\qquad 0\le v_m(x)<M_m,
\]
并且该 $x$ 必为 $\Fold_m(\omega)$。等价地，
\[
v_m(\Fold_m(\omega))\equiv N(\omega)\pmod{M_m},
\qquad
0\le v_m(\Fold_m(\omega))<M_m,
\qquad
c=\mathbf 1_{\{N(\omega)\ge M_m\}}.
\]
\leanverified{stableValue\_injective}
\leanverified{stableAdd\_value\_mod}
\leanverified{paper\_fold\_zeckendorf\_mod\_decomposition\_unique}
\end{theorem}

\begin{proof}
由 $0\le N(\omega)\le T_m=F_{m+3}-2<M_m+F_{m+1}\le 2M_m$，可知
$c:=\mathbf 1_{\{N(\omega)\ge M_m\}}\in\{0,1\}$ 唯一确定。令 $r:=N(\omega)-cM_m$，则 $0\le r<M_m$。
在区间 $[0,M_m)$ 上，Zeckendorf 表示不会使用权重 $F_{m+2}$，因此存在唯一 $x\in X_m$ 使 $v_m(x)=r$。
另一方面，$Z(N(\omega))$ 的前 $m$ 位正是 $\Fold_m(\omega)$；当且仅当 $c=c_{m+1}=1$ 时出现顶位 $F_{m+2}$，截断后余数即为 $r$，从而 $v_m(\Fold_m(\omega))=r$。唯一性由 $c$ 与 Zeckendorf 表示唯一性推出。
\end{proof}

\begin{proposition}[补位对称诱导的 Fourier 相位刚性]\label{prop:fold-fiber-fourier-phase-rigidity}
沿用命题 \ref{prop:fold-fiber-count-reciprocity} 的记号，并令
\[
\alpha_m:=F_{m+1}-2\in\ZZ/M_m\ZZ.
\]
将 \(f_m\) 视为加法群 \(\ZZ/M_m\ZZ\) 上的类函数，并定义其离散 Fourier 变换
\[
\widehat f_m(k):=\sum_{r\in\ZZ/M_m\ZZ} f_m(r)\exp\!\left(-\frac{2\pi i kr}{M_m}\right),
\qquad k\in\ZZ/M_m\ZZ.
\]
则对任意 \(k\in\ZZ/M_m\ZZ\) 有严格恒等式
\[
\boxed{\;
\widehat f_m(k)
=\exp\!\left(-\frac{2\pi i k\alpha_m}{M_m}\right)\overline{\widehat f_m(k)}\;}
\]
从而
\[
\boxed{\;
\exp\!\left(\frac{\pi i k\alpha_m}{M_m}\right)\widehat f_m(k)\in\RR\;}
\]
即每个 Fourier 系数的相位被锁定在一条固定的实直线子空间中。
\end{proposition}
\begin{proof}
由命题 \ref{prop:fold-fiber-count-reciprocity} 的反射对称性 \(f_m(r)=f_m(\alpha_m-r)\)，有
\[
\widehat f_m(k)
=\sum_{r} f_m(\alpha_m-r)\exp\!\left(-\frac{2\pi i kr}{M_m}\right)
=\exp\!\left(-\frac{2\pi i k\alpha_m}{M_m}\right)\sum_{r} f_m(r)\exp\!\left(\frac{2\pi i kr}{M_m}\right).
\]
最后一项为 \(\overline{\widehat f_m(k)}\)，故得第一式；第二式由移项并乘以半相位因子得到。证毕。
\end{proof}
\leanverified{paper\_fold\_fiber\_fourier\_phase\_rigidity}

\begin{corollary}[模余类刚性与有限分辨率谓词的周期性]\label{cor:fold-mod-rigidity-periodic-predicates}
固定 $m\ge 1$，定义剩余映射
\[
r_m:\Omega_m\to \ZZ/M_m\ZZ,\qquad r_m(\omega):=N(\omega)\bmod M_m.
\]
令 \(s_m:\ZZ/M_m\ZZ\to X_m\) 为如下规范截面：对任意剩余类 \(\bar r\)，取其唯一代表元 \(r\in\{0,1,\dots,M_m-1\}\)，并令 \(s_m(\bar r)\) 为满足 \(v_m(x)=r\) 的唯一 \(x\in X_m\)。
则
\[
\Fold_m\ =\ s_m\circ r_m.
\]
因此对任意集合 \(Y\) 与任意映射 \(g:X_m\to Y\)，复合 \(g\circ \Fold_m\) 仅依赖于 \(r_m(\omega)\)。
特别地，对任意谓词 \(P\subseteq X_m\)，存在余类子集 \(R_P\subseteq \ZZ/M_m\ZZ\) 使得
\[
\Fold_m^{-1}(P)=r_m^{-1}(R_P).
\]
从而把 \(\ind_{\{\Fold_m(\omega)\in P\}}\) 视为 \(N(\omega)\) 的函数时，其必为周期 \(M_m\) 的周期函数。
\end{corollary}

\begin{proof}
由定理 \ref{thm:fold-zeckendorf-mod-decomposition-unique}，对任意 \(\omega\in\Omega_m\) 有
\(
v_m(\Fold_m(\omega))\equiv N(\omega)\ (\mathrm{mod}\ M_m)
\)
且 \(v_m(\Fold_m(\omega))\in\{0,\dots,M_m-1\}\)。
而 \(s_m\) 恰按该余数选择唯一的 \(X_m\) 元素，因此 \(\Fold_m(\omega)=s_m(r_m(\omega))\)。
其余断言由因子化直接推出。证毕。
\leanverified{fold\_eq\_of\_weight\_mod\_eq}
\leanverified{paper\_fold\_mod\_rigidity\_periodic\_predicates}
\end{proof}

\begin{proposition}[端点 $M_m-1$ 的子集和表示唯一性]\label{prop:fold-endpoint-Mm-minus-one-unique}
对任意 $m\ge 2$，方程
\[
\sum_{k=1}^{m}\omega_kF_{k+1}=M_m-1=F_{m+2}-1,\qquad \omega\in\Omega_m
\]
恰有一个解，并由交错模式给出
\[
(\omega_m,\omega_{m-1},\dots,\omega_1)=
\begin{cases}
(1,0,1,0,\dots,1,0), & m\ \text{为偶数},\\
(1,0,1,0,\dots,0,1), & m\ \text{为奇数}.
\end{cases}
\]
等价地，
\[
F_{m+2}-1=F_{m+1}+F_{m-1}+F_{m-3}+\cdots.
\]
\end{proposition}

\begin{proof}
设 $S=\sum_{k=1}^{m}\omega_kF_{k+1}=F_{m+2}-1$。注意 $\sum_{k=1}^{m-1}F_{k+1}=F_{m+2}-2$，故若 $\omega_m=0$ 则 $S\le F_{m+2}-2$ 矛盾，因而 $\omega_m=1$。
于是
\[
\sum_{k=1}^{m-1}\omega_kF_{k+1}=S-F_{m+1}=F_m-1.
\]
同理 $\sum_{k=1}^{m-2}F_{k+1}=F_m-2$，故 $\omega_{m-1}=1$ 不可能，从而 $\omega_{m-1}=0$。继续对剩余等式递归同样论证，得到 $\omega_{m-2}=1,\omega_{m-3}=0,\dots$ 的强制交错结构，因而解唯一，并给出所述展开。
\leanverified{alternating\_fib\_sum\_even}
\leanverified{alternating\_fib\_sum\_odd}
\leanverified{paper\_fold\_endpoint\_seeds}
\leanverified{paper\_fold\_endpoint\_Mm\_minus\_one\_unique\_seeds}
\leanverified{paper\_fold\_endpoint\_Mm\_minus\_one\_unique\_package}
\leanverified{paper\_fold\_endpoint\_Mm\_minus\_one\_unique}
\leanverified{paper\_fold\_endpoint\_mm\_minus\_one\_unique}
\end{proof}

\begin{theorem}[顶两位 Zeckendorf 统计的三分律与有限修正]\label{thm:fold-top-two-zeckendorf-trisect}
在 $\Omega_m$ 上取均匀测度，并令
\[
A_m:=\#\{\omega\in\Omega_m:\ N(\omega)\ge M_m\}.
\]
则对一切 $m\ge 2$，
\[
\#\{(\zeta_m,\zeta_{m+1})=(0,1)\}=A_m,\qquad
\#\{(\zeta_m,\zeta_{m+1})=(1,0)\}=2^m-2A_m-1,\qquad
\#\{(\zeta_m,\zeta_{m+1})=(0,0)\}=A_m+1,
\]
且
\[
A_m=
\begin{cases}
\dfrac{4^{m/2}-1}{3}, & m\ \text{为偶数},\\[6pt]
\dfrac{2\bigl(4^{(m-1)/2}-1\bigr)}{3}, & m\ \text{为奇数}.
\end{cases}
\]
因此极限分布满足
\[
\lim_{m\to\infty}\PP((\zeta_m,\zeta_{m+1})=(0,0))
=\lim_{m\to\infty}\PP((\zeta_m,\zeta_{m+1})=(1,0))
=\lim_{m\to\infty}\PP((\zeta_m,\zeta_{m+1})=(0,1))=\frac13,
\]
并且有限尺寸修正为
\[
\PP((\zeta_m,\zeta_{m+1})=(0,1))=
\begin{cases}
\dfrac{1-2^{-m}}{3}, & m\ \text{为偶数},\\[6pt]
\dfrac{1-2^{-(m-1)}}{3}, & m\ \text{为奇数},
\end{cases}
\]
\[
\PP((\zeta_m,\zeta_{m+1})=(1,0))=
\begin{cases}
\dfrac{1-2^{-m}}{3}, & m\ \text{为偶数},\\[6pt]
\dfrac{1+2^{-m}}{3}, & m\ \text{为奇数},
\end{cases}
\]
\[
\PP((\zeta_m,\zeta_{m+1})=(0,0))=
\begin{cases}
\dfrac{1+2^{-(m-1)}}{3}, & m\ \text{为偶数},\\[6pt]
\dfrac{1+2^{-m}}{3}, & m\ \text{为奇数}.
\end{cases}
\]
\leanverified{hiddenBitCount\_even\_closed}
\leanverified{hiddenBitCount\_odd\_closed}
\leanverified{complement\_hiddenBitCount}
\leanverified{hiddenBitCount\_lt\_pow}
\leanverified{hiddenBitCount\_pos}
\leanverified{paper\_fold\_top\_two\_zeckendorf\_trisect}
\end{theorem}

\begin{proof}
先给出 $A_m$ 的递推。若 $\omega_m=0$，则 $N(\omega)\le \sum_{k=1}^{m-1}F_{k+1}=F_{m+2}-2<M_m$，故事件 $\{N(\omega)\ge M_m\}$ 蕴含 $\omega_m=1$。于是
\[
A_m=\#\Bigl\{\omega'\in\{0,1\}^{m-1}:\ \sum_{k=1}^{m-1}\omega'_kF_{k+1}\ge F_m\Bigr\}.
\]
取补集并对 $\{0,1\}^{m-1}$ 施行逐位取反对合 $\varepsilon\mapsto \bar\varepsilon$，由
$\sum_{k=1}^{m-1}F_{k+1}=F_{m+2}-2$ 得
\[
A_m=\#\Bigl\{\varepsilon\in\{0,1\}^{m-1}:\ \sum_{k=1}^{m-1}\varepsilon_kF_{k+1}\le F_{m+1}-2\Bigr\}.
\]
按 $\varepsilon_{m-1}$ 分解：若 $\varepsilon_{m-1}=0$，则前 $m-2$ 位任意皆满足不等式（因 $\sum_{k=1}^{m-2}F_{k+1}=F_{m+1}-2$），贡献 $2^{m-2}$；
若 $\varepsilon_{m-1}=1$，则剩余需满足
\[
\sum_{k=1}^{m-2}\varepsilon_kF_{k+1}\le (F_{m+1}-2)-F_m=F_{m-1}-2,
\]
其计数恰为 $A_{m-2}$。因此得到递推
\[
A_m=2^{m-2}+A_{m-2}\qquad(m\ge 4),
\]
并由直接核算 $A_2=1,A_3=2$ 解得所示闭式。

其次把 $A_m$ 与顶两位联系。由 Zeckendorf 约束 $c_kc_{k+1}=0$，$(\zeta_m,\zeta_{m+1})$ 仅有三态，并且按权重区间划分有
\[
(\zeta_m,\zeta_{m+1})=(0,1)\iff N(\omega)\ge F_{m+2}=M_m,
\]
\[
(\zeta_m,\zeta_{m+1})=(1,0)\iff F_{m+1}\le N(\omega)<F_{m+2}=M_m,
\qquad
(\zeta_m,\zeta_{m+1})=(0,0)\iff N(\omega)<F_{m+1}.
\]
第一式给出 $\#\{(\zeta_m,\zeta_{m+1})=(0,1)\}=A_m$。由逐位取反对合 $\omega\mapsto \bar\omega$ 满足
$N(\bar\omega)=T_m-N(\omega)$，可得
\[
\#\{N(\omega)<F_{m+1}\}=\#\{N(\omega)\ge T_m-(F_{m+1}-1)+1\}=\#\{N(\omega)\ge F_{m+2}-1\}.
\]
右端等于 $A_m+\#\{N(\omega)=F_{m+2}-1\}$，而 $\#\{N(\omega)=F_{m+2}-1\}=1$ 由命题 \ref{prop:fold-endpoint-Mm-minus-one-unique} 给出，从而
\(
\#\{(\zeta_m,\zeta_{m+1})=(0,0)\}=A_m+1
\)。
再由总数 $2^m$ 得到
\[
\#\{(\zeta_m,\zeta_{m+1})=(1,0)\}=2^m-\bigl(A_m+(A_m+1)\bigr)=2^m-2A_m-1.
\]
将 $A_m$ 的闭式代入即可得到两条概率修正公式。
\end{proof}

\begin{corollary}[折叠溢出率与缺口期望的闭式]\label{cor:fold-overflow-gap-rate-expectation}
在 $\Omega_m$ 的均匀测度下，令
\[
\Delta_m(\omega):=N(\omega)-v_m(\Fold_m(\omega))\in\{0,M_m\}
\]
为折叠缺口，则对一切 $m\ge 2$，
\[
\PP(\Delta_m=M_m)=\frac{A_m}{2^m}=
\begin{cases}
\dfrac{1-2^{-m}}{3}, & m\ \text{为偶数},\\[6pt]
\dfrac{1-2^{-(m-1)}}{3}, & m\ \text{为奇数},
\end{cases}
\qquad
\EE[\Delta_m]=M_m\,\PP(\Delta_m=M_m).
\]
并且
\[
\PP\bigl(v_m(\Fold_m(\omega))<F_{m+1}\bigr)=
\begin{cases}
\dfrac{1+2^{-(m-1)}}{3}, & m\ \text{为偶数},\\[6pt]
\dfrac{1+2^{-m}}{3}, & m\ \text{为奇数}.
\end{cases}
\]
\[
\PP\bigl(F_{m+1}\le N(\omega)<M_m\bigr)=
\begin{cases}
\dfrac{1-2^{-m}}{3}, & m\ \text{为偶数},\\[6pt]
\dfrac{1+2^{-m}}{3}, & m\ \text{为奇数}.
\end{cases}
\]
\end{corollary}

\begin{proof}
由定理 \ref{thm:fold-zeckendorf-mod-decomposition-unique}，$\Delta_m=cM_m$ 且 $\Delta_m=M_m\iff N(\omega)\ge M_m$，故 $\PP(\Delta_m=M_m)=A_m/2^m$。代入定理 \ref{thm:fold-top-two-zeckendorf-trisect} 的 $A_m$ 化简得到所示闭式，期望等式由 $\Delta_m\in\{0,M_m\}$ 立得。第二个概率等式等价于事件 $(\zeta_m,\zeta_{m+1})=(0,0)$，第三个概率等式等价于事件 $(\zeta_m,\zeta_{m+1})=(1,0)$，均由定理 \ref{thm:fold-top-two-zeckendorf-trisect} 给出。
\leanverified{three\_mul\_hiddenBitCount\_add}
\leanverified{paper\_fold\_overflow\_gap\_rate\_expectation}
\end{proof}

\begin{proposition}[零剩余纤维的线性退化]\label{prop:fold-zero-fiber-linear}
令
\[
f_m(0):=\#\{\omega\in\Omega_m:\ \Fold_m(\omega)=0^m\},
\]
则对任意 $m\ge 2$，
\[
f_m(0)=1+\Bigl\lfloor\frac{m}{2}\Bigr\rfloor.
\]
\end{proposition}

\begin{proof}
由定理 \ref{thm:fold-zeckendorf-mod-decomposition-unique}，$\Fold_m(\omega)=0^m$ 当且仅当 $v_m(\Fold_m(\omega))=0$，等价于 $N(\omega)\equiv 0\pmod{M_m}$。
又因 $0\le N(\omega)<2M_m$，上述同余等价于 $N(\omega)\in\{0,M_m\}$。其中 $N(\omega)=0$ 仅由全零 $\omega=0^m$ 给出，贡献 $1$。
剩余需计数
\[
B_m:=\#\{\omega\in\Omega_m:\ N(\omega)=M_m=F_{m+2}\}.
\]
若 $\omega_m=0$ 则 $N(\omega)\le F_{m+2}-2$ 矛盾，故必有 $\omega_m=1$，从而
\[
\sum_{k=1}^{m-1}\omega_kF_{k+1}=F_{m+2}-F_{m+1}=F_m.
\]
在右侧表示中，若 $\omega_{m-1}=1$，则左侧含项 $F_m$，其余位必须为 $0$，得到一种表示；若 $\omega_{m-1}=0$，则需用 $\{F_2,\dots,F_{m-1}\}$ 表示 $F_m$，这与参数 $m-2$ 的同类计数等价，故
\[
B_m=1+B_{m-2}\qquad(m\ge 4),
\]
并由 $B_2=1,B_3=1$ 解得 $B_m=\lfloor m/2\rfloor$。因此 $f_m(0)=1+B_m$。
\leanverified{paper\_fold\_zero\_fiber\_linear}
\leanverified{paper\_fold\_zero\_fiber\_linear\_extended}
\leanverified{weightSumAtMm}
\leanverified{weightSumAtMm\_two}
\leanverified{weightSumAtMm\_three}
\leanverified{weightSumAtMm\_four}
\leanverified{weightSumAtMm\_five}
\leanverified{weightSumAtMm\_eq\_div\_two\_small}
\end{proof}

\begin{proposition}[纤维计数的互易律与反射对称性]\label{prop:fold-fiber-count-reciprocity}
\leanverified{paper\_fold\_fiber\_count\_reciprocity}
对任意 $m\ge 2$ 与任意 $r\in\{0,1,\dots,M_m-1\}$，令
\[
f_m(r):=\#\{\omega\in\Omega_m:\ v_m(\Fold_m(\omega))=r\},
\qquad
a_m(s):=\#\{\omega\in\Omega_m:\ N(\omega)=s\}\quad(0\le s\le T_m),
\]
则成立互易恒等式
\[
f_m(r)=a_m(r)+a_m(M_m+r),
\]
并且
\[
f_m(r)=a_m(r)+a_m(F_{m+1}-2-r),
\qquad
f_m(r)=f_m(F_{m+1}-2-r).
\]
\leanpartial{weightCongruenceCount\_complement}{条件 m≥2，r≤F\_{m+1}-2；互易恒等式 f\_m(r)=a\_m(r)+a\_m(M\_m+r) 未形式化}
\leanverified{fiberMultiplicity\_complement}
\leanverified{stableValue\_Fold\_add\_complement}
\leanverified{complementAction}
\leanverified{complementAction\_involutive}
\leanverified{fiberMultiplicity\_complementAction}
\leanverified{fiberMultiplicity\_complement\_parity}
\leanverified{popcount\_complement}
\leanverified{paper\_pom\_fiber\_reciprocity\_package}
\leanverified{paper\_pom\_fiber\_histogram\_palindrome}
\leanverified{paper\_folding\_weight\_count\_palindrome}
\end{proposition}

% --- Distillation writeback ---
\begin{definition}[fold-aware 跨分辨率限制]\label{distill:wang_zahl-sticky_reduction_and_scale_saturated_closure-fold-aware-cross-resolution}
\textbf{结果 17；日期：2026-04-23T03:31:39+08:00。依赖状态：construction.}
设 \(m\ge h\ge 1\)，并令 \(Z(N(\omega))=(c_1,c_2,\ldots)\) 为 \(N(\omega)\) 的 Zeckendorf 数字序列。定义 fold-aware 跨分辨率限制
\[
\Fold^{\rm aw}_{m\to h}:\Omega_m\longrightarrow X_h,
\qquad
\Fold^{\rm aw}_{m\to h}(\omega):=(c_1,\ldots,c_h).
\]
同时记稳定前缀投影为
\[
\pr_{m,h}:X_m\longrightarrow X_h,
\qquad
\pr_{m,h}(x_1,\ldots,x_m):=(x_1,\ldots,x_h).
\]
该对象不是把原始输入字 \(\omega\) 截成前 \(h\) 位后再施行 \(\Fold_h\)；它先在同一个整数 \(N(\omega)\) 上取 Zeckendorf 稳定展开，再取稳定前缀。因此它避免了把朴素输入截断误认为跨分辨率不变量的幻觉。
\end{definition}

\begin{proposition}[fold-aware 限制与稳定前缀交换]\label{distill:wang_zahl-sticky_reduction_and_scale_saturated_closure-fold-aware-prefix-commutation}
\textbf{结果 18；日期：2026-04-23T03:31:39+08:00。依赖状态：rigidity.}
对任意 \(m\ge h\ge 1\)，有严格交换关系
\[
\Fold^{\rm aw}_{m\to h}=\pr_{m,h}\circ\Fold_m.
\]
因此，若 \(\gamma_h:X_h\to Y_h\) 是任意 coarse carrier 映射，则
\[
\gamma_h\circ\Fold^{\rm aw}_{m\to h}
=
(\gamma_h\circ\pr_{m,h})\circ\Fold_m.
\]
特别地，所有跨分辨率 sticky/non-sticky 测试若经由 \(\gamma_h\) 读出，都可完全下降到稳定代表 \(\Fold_m(\omega)\) 的前缀，而不依赖原始微态字的朴素前缀。
\end{proposition}

% --- Distillation writeback ---
\begin{proposition}[投影 fold congruence 的尾部 lift 审计]\label{prop:distill-euclid-elements-fold-prefix-tail-lift-audit}
\textbf{Dependency status: rigidity.}
设 \(m\ge h\ge1\)。沿用定义 \ref{distill:wang_zahl-sticky_reduction_and_scale_saturated_closure-fold-aware-cross-resolution} 的
\(\Fold^{\rm aw}_{m\to h}\) 与稳定前缀投影 \(\pr_{m,h}\)。再定义尾部读出
\[
\operatorname{Tail}_{m,h}:X_m\longrightarrow \{0,1\}^{m-h},
\qquad
\operatorname{Tail}_{m,h}(x_1,\ldots,x_m):=(x_{h+1},\ldots,x_m),
\]
其中 \(h=m\) 时右端理解为单点集合。则对任意 \(\omega,\omega'\in\Omega_m\)，有
\[
\Fold_m(\omega)=\Fold_m(\omega')
\]
当且仅当同时满足
\[
\Fold^{\rm aw}_{m\to h}(\omega)=\Fold^{\rm aw}_{m\to h}(\omega')
\]
与
\[
\operatorname{Tail}_{m,h}(\Fold_m(\omega))
=
\operatorname{Tail}_{m,h}(\Fold_m(\omega')).
\]
因此，深度 \(h\) 的 projected congruence 只证明稳定前缀相等；它成为深度 \(m\) 的 lifted equality，当且仅当尾部 residual 同时消失。
\end{proposition}

% --- Distillation writeback ---
% Distillation timestamp: 2026-04-23T12:04:19+08:00; pipeline index: 35.
\begin{theorem}[投影折叠同余的最小尾部严格化]\label{distill:euclid_elements-lifted_congruence_and_residual_rigidity_families-fold-projected-congruence-minimal-tail-strictification}
\textbf{Dependency status: rigidity.}
设 \(m\ge h\ge 1\)，并取非空有限族 \(E\subseteq\Omega_m\)。沿用
\(\Fold^{\rm aw}_{m\to h}\)、\(\pr_{m,h}\) 与
\(\operatorname{Tail}_{m,h}\) 的记号。对每个 \(p\in X_h\)，定义投影纤维中的尾部残余集合
\[
\mathcal T_E(p)
:=
\left\{
\operatorname{Tail}_{m,h}\bigl(\Fold_m(\omega)\bigr):
\omega\in E,\ \Fold^{\rm aw}_{m\to h}(\omega)=p
\right\},
\]
并令
\[
\mu_{m,h}(E):=\max_{p\in X_h}\bigl|\mathcal T_E(p)\bigr|.
\]
称一个有限寄存器 \(\rho:E\to R\) 为该投影同余的尾部 lift 审计，若对任意
\(\omega,\omega'\in E\) 都有
\[
\Fold^{\rm aw}_{m\to h}(\omega)=\Fold^{\rm aw}_{m\to h}(\omega'),\qquad
\rho(\omega)=\rho(\omega')
\quad\Longrightarrow\quad
\Fold_m(\omega)=\Fold_m(\omega').
\]
则：
\begin{enumerate}
  \item 任意尾部 lift 审计都满足
  \[
  |R|\ge \mu_{m,h}(E),
  \]
  且存在达到等号的尾部 lift 审计；
  \item 深度 \(h\) 的 projected congruence 在 \(E\) 上已经推出深度 \(m\) 的 lifted equality，当且仅当
  \[
  \mu_{m,h}(E)=1;
  \]
  \item 若 \(\mu_{m,h}(E)>1\)，则任意使该投影同余失败的坏对
  \[
  \Fold^{\rm aw}_{m\to h}(\omega)=\Fold^{\rm aw}_{m\to h}(\omega'),\qquad
  \Fold_m(\omega)\ne \Fold_m(\omega')
  \]
  必且只必携带非零尾部 residual：
  \[
  \operatorname{Tail}_{m,h}\bigl(\Fold_m(\omega)\bigr)
e
  \operatorname{Tail}_{m,h}\bigl(\Fold_m(\omega')\bigr).
  \]
  因而任何仅经由 \(\Fold^{\rm aw}_{m\to h}\) 因子化的 coarse readout 都不能区分这样的坏对。
\end{enumerate}
\end{theorem}
\begin{proof}
固定 \(p\in X_h\)。若 \(\tau,\tau'\in\mathcal T_E(p)\) 且 \(\tau\ne\tau'\)，则可取
\(\omega,\omega'\in E\) 使
\[
\Fold^{\rm aw}_{m\to h}(\omega)=
\Fold^{\rm aw}_{m\to h}(\omega')=p,
\]
且两者尾部读出分别为 \(\tau,\tau'\)。由命题
\ref{distill:wang_zahl-sticky_reduction_and_scale_saturated_closure-fold-aware-prefix-commutation}，投影同余正是稳定前缀相等；再由命题
\ref{prop:distill-euclid-elements-fold-prefix-tail-lift-audit}，在同一前缀下，深度 \(m\) 的折叠值相等当且仅当尾部也相等。因此上述两点的 lifted equality 失败。若 \(\rho\) 是 lift 审计，则它必须给这些不同尾部赋予不同寄存器值。于是
\(|R|\ge |\mathcal T_E(p)|\)。对所有 \(p\) 取最大值得到
\(|R|\ge\mu_{m,h}(E)\)。

反过来，取一个集合 \(R\) 满足 \(|R|=\mu_{m,h}(E)\)。对每个 \(p\) 任选单射
\[
\iota_p:\mathcal T_E(p)\hookrightarrow R,
\]
并定义
\[
\rho(\omega):=\iota_{\Fold^{\rm aw}_{m\to h}(\omega)}
\left(
\operatorname{Tail}_{m,h}(\Fold_m(\omega))
\right).
\]
若两点具有相同 projected congruence 且 \(\rho\) 值相同，则它们位于同一个 \(p\)-纤维内，并因 \(\iota_p\) 单射而具有相同尾部。命题
\ref{prop:distill-euclid-elements-fold-prefix-tail-lift-audit} 遂给出 \(\Fold_m(\omega)=\Fold_m(\omega')\)。故等号可达。

第二条是第一条的退化情形。若 \(\mu_{m,h}(E)=1\)，则每个投影前缀下只有一个尾部 residual，故 projected congruence 自动给出前缀与尾部同时相等，从而给出 lifted equality。若 \(\mu_{m,h}(E)>1\)，则存在同一 \(p\) 下两个不同尾部，取相应代表即得到 projected congruent 但不 lifted equal 的坏对。

第三条再次由 \(X_m\subseteq\{0,1\}^m\) 中一个词由其前缀与尾部唯一决定得到：在前缀已经相等的前提下，折叠值不等价于尾部不等。若 \(g:E\to Y\) 仅经由 \(\Fold^{\rm aw}_{m\to h}\) 因子化，则存在 \(\bar g:X_h\to Y\) 使
\(g=\bar g\circ\Fold^{\rm aw}_{m\to h}\)，故上述坏对必有相同 \(g\)-值。这正是投影同余不能无审计转移为规范相等的尾部 residual 机制。证毕。
\end{proof}

% --- End distillation writeback ---

\begin{proof}
由命题 \ref{distill:wang_zahl-sticky_reduction_and_scale_saturated_closure-fold-aware-prefix-commutation}，
\[
\Fold^{\rm aw}_{m\to h}=\pr_{m,h}\circ\Fold_m.
\]
于是 projected congruence 等价于
\[
\pr_{m,h}(\Fold_m(\omega))=\pr_{m,h}(\Fold_m(\omega')).
\]
另一方面，任意 \(x\in X_m\subseteq\{0,1\}^m\) 由其前缀 \(\pr_{m,h}(x)\) 与尾部 \(\operatorname{Tail}_{m,h}(x)\) 唯一决定。因此两个折叠值 \(\Fold_m(\omega)\) 与 \(\Fold_m(\omega')\) 相等，当且仅当前缀和尾部均相等。证毕。
\end{proof}

% --- End distillation writeback ---

\begin{proof}
由定义，\(\Fold_m(\omega)\) 正是 \(Z(N(\omega))\) 的前 \(m\) 个 Zeckendorf 稳定位，即
\[
\Fold_m(\omega)=(c_1,\ldots,c_m).
\]
对其施加稳定前缀投影 \(\pr_{m,h}\)，得到
\[
\pr_{m,h}(\Fold_m(\omega))=(c_1,\ldots,c_h).
\]
右端按定义正是 \(\Fold^{\rm aw}_{m\to h}(\omega)\)。故第一式成立。第二式仅是在第一式两端复合 \(\gamma_h\)。最后一句说明：若 carrier 读出为 \(\gamma_h\)，则其跨尺度值已经由稳定前缀决定；这里未使用、也不需要使用原始输入字的前 \(h\) 位，因此不会把 \(\Fold_h(\omega_1,\ldots,\omega_h)\) 误作跨分辨率不变量。证毕。
\end{proof}

% --- End distillation writeback ---

\begin{proof}
由定理 \ref{thm:fold-zeckendorf-mod-decomposition-unique}，$v_m(\Fold_m(\omega))$ 是 $N(\omega)$ 模 $M_m$ 的规范代表，且 $0\le N(\omega)<2M_m$，故事件 $\{v_m(\Fold_m(\omega))=r\}$ 等价于 $\{N(\omega)\in\{r,M_m+r\}\}$，从而得到第一式。
对合 $\omega\mapsto\bar\omega$ 满足 $N(\bar\omega)=T_m-N(\omega)$，故 $a_m(s)=a_m(T_m-s)$。注意
\[
T_m-(M_m+r)=(F_{m+3}-2)-(F_{m+2}+r)=F_{m+1}-2-r,
\]
代入 $a_m(M_m+r)=a_m(F_{m+1}-2-r)$ 即得第二式。将第二式在 $r$ 与 $F_{m+1}-2-r$ 间交换即可推出反射对称性。
\end{proof}


\subsubsection{折叠的局部重写实现：终止、合流与顺序无关}\label{subsec:fold-local-rewrite}
定义 \ref{def:fold-word} 以“先取值、再取 Zeckendorf 规范形”给出 $\Fold_m$。为将其落实为\textbf{局部、可执行}的归一化过程，并澄清“归约顺序不影响正规形”，我们把输入窗口嵌入到允许中间态出现冗余数字的配置空间上实施值保持重写。

这一实现顺序也很适合主论文的整体叙述：先说明对象可由哪些局部许可重写得到，再说明为何不同归约顺序会坍缩到同一正规形。前者是 constructibility，后者是 rigidity；两层若不拆开，折叠就会被误读成“定义一个商映射”，而不是“通过合法局部作图获得规范对象”。

\paragraph{配置空间与值函数}
令
\[
\mathcal{D}:=\Bigl\{a=(a_1,a_2,\dots): a_k\in\NN,\ \text{仅有限个 }a_k\neq 0\Bigr\}
\]
为有限支撑的非负整数数字配置。定义值同态
\[
\mathrm{Val}(a):=\sum_{k\ge 1} a_k F_{k+1}.
\]
对给定 $\omega\in\Omega_m$，记其嵌入为
\[
\iota_m(\omega):=(\omega_1,\dots,\omega_m,0,0,\dots)\in\mathcal{D},
\]
则 $\mathrm{Val}(\iota_m(\omega))=N(\omega)$。

\paragraph{局部重写规则（值保持）}
对 $a\in\mathcal{D}$，允许在满足条件时对局部模式做如下替换（其中 $e_k$ 为第 $k$ 位单位向量）：
\begin{enumerate}
  \renewcommand{\theenumi}{R\arabic{enumi}}
  \renewcommand{\labelenumi}{(\theenumi)}
  \item \label{rule:fold:adj}\textbf{相邻合并：}若 $a_k\ge 1$ 且 $a_{k+1}\ge 1$，则
  \[
  a\ \to\ a-e_k-e_{k+1}+e_{k+2}.
  \]
  \item \label{rule:fold:dedup1}\textbf{底位去重：}若 $a_1\ge 2$，则
  \[
  a\ \to\ a-2e_1+e_2.
  \]
  \item \label{rule:fold:dedup2}\textbf{第二位去重：}若 $a_2\ge 2$，则
  \[
  a\ \to\ a-2e_2+e_1+e_3.
  \]
  \item \label{rule:fold:dedupk}\textbf{一般位去重（$k\ge 3$）：}若 $a_k\ge 2$，则
  \[
  a\ \to\ a-2e_k+e_{k-2}+e_{k+1}.
  \]
\end{enumerate}
每条规则都只改动有限个坐标，因此是严格的\textbf{局部}重写；并且它们分别对应恒等式
\(
F_{k+1}+F_{k+2}=F_{k+3},
2F_2=F_3,
2F_3=F_2+F_4,
2F_{k+1}=F_{k-1}+F_{k+2}
\)，因此每一步都保持 $\mathrm{Val}$ 不变。

\paragraph{不可约项与 Zeckendorf 规范形}
记 $X_\infty=\{(c_1,c_2,\dots)\in\{0,1\}^{\mathbb{N}}:\ c_kc_{k+1}=0\}$。上述规则恰在出现相邻 $11$ 或某位数字 $\ge 2$ 时可用，因此不可约项当且仅当属于 $X_\infty$。

\begin{proposition}[分辨率稳定类型的逆极限本体]\label{prop:fold-inverse-limit-xm-xinfty}
对每个 \(m\ge 1\) 令 \(r_{m+1\to m}:X_{m+1}\to X_m\) 为前缀截断 \(r_{m+1\to m}(a_1\cdots a_{m+1})=a_1\cdots a_m\)。
则 \((X_m,r_{m+1\to m})\) 构成逆系统，并存在自然同构
\[
\varprojlim_m X_m\ \cong\ X_\infty,
\]
其中 \(X_\infty\to X_m\) 的结构映射为取前 \(m\) 位前缀。
此外，对任意 \(m\ge 1\)，Zeckendorf 截断映射
\(
N\mapsto \pi_m(Z(N))
\)
在区间 \(\{0,1,\dots,F_{m+2}-1\}\) 上与 \(X_m\) 构成双射，其逆映射为 \(w\mapsto V_m(w)\)（\S\ref{subsec:emergent-arithmetic-zeckendorf-value-address}；并见引理 \ref{lem:zeckendorf-range-bijection}）。
\leanverified{paper\_fold\_inverse\_limit\_xm\_xinfty}
\leanverified{prefixWord\_stableValue\_coherent}
\leanverified{prefixWord\_surjective}
\end{proposition}
\begin{proof}
由于“无相邻 \(11\)”对前缀封闭，\(r_{m+1\to m}\) 良定义，从而得到逆系统。
给定相容族 \((w_m)_{m\ge 1}\in\varprojlim_m X_m\)，对每个坐标 \(i\ge 1\) 令 \(a_i\) 为任意 \(m\ge i\) 时 \(w_m\) 的第 \(i\) 位。
相容性保证该定义与 \(m\) 的选择无关，因此得到唯一的无限词 \(a=(a_i)_{i\ge 1}\in\{0,1\}^{\NN}\)；
又因每个 \(w_m\) 无相邻 \(11\)，故 \(a\in X_\infty\)。
反之，对任意 \(a\in X_\infty\)，其前缀 \(\pi_m(a)\in X_m\) 形成相容族 \((\pi_m(a))_m\in\varprojlim X_m\)。
两构造互为逆，得到同构。
最后一条双射性由 \(V_m\) 的范围双射与 Zeckendorf 唯一性给出（引理 \ref{lem:zeckendorf-range-bijection}）。证毕。
\end{proof}

\begin{proposition}[一致纤维集的前缀覆盖数与外置比特率上界]\label{prop:fold-consistent-fiber-hausdorff-prefix-count}
令 \(\Omega_\infty:=\{0,1\}^{\NN}\)，并以 \(\tau_{\infty\to m}:\Omega_\infty\to\Omega_m\) 表示取前 \(m\) 位前缀。
给定 \(x_\infty\in X_\infty\cong\varprojlim X_m\)，记 \(x_m:=\pi_m(x_\infty)\in X_m\)。
定义与 \(x_\infty\) 相容的一致纤维集
\[
\mathcal F(x_\infty)
:=\Bigl\{\omega\in\Omega_\infty:\ \Fold_m\bigl(\tau_{\infty\to m}(\omega)\bigr)=x_m\ \text{对所有 }m\ge 1\Bigr\}.
\]
对每个 \(m\ge 1\) 令
\[
N_m(x_\infty)
:=\#\Bigl\{a\in\Omega_m:\ [a]\cap \mathcal F(x_\infty)\neq\varnothing\Bigr\}
=\#\Bigl\{\tau_{\infty\to m}(\omega):\ \omega\in\mathcal F(x_\infty)\Bigr\},
\]
其中 \([a]\subset\Omega_\infty\) 为柱集。
则有纯计数上界
\[
N_m(x_\infty)\ \le\ d_m(x_m):=\bigl|\Fold_m^{-1}(x_m)\bigr|.
\]
并且对任意 \(\omega\in\mathcal F(x_\infty)\) 都存在确定性解码使
\[
K\bigl(\tau_{\infty\to m}(\omega)\mid x_\infty,m\bigr)
\ \le\
\log_2 d_m(x_m)+O(\log m).
\]
因此
\[
\liminf_{m\to\infty}\frac{K(\tau_{\infty\to m}(\omega)\mid x_\infty,m)}{m}
\ \le\
\liminf_{m\to\infty}\frac{\log_2 d_m(x_m)}{m}.
\]
\leanverified{paper\_fold\_consistent\_fiber\_hausdorff\_prefix\_count}
\end{proposition}
\begin{proof}
由定义，若 \([a]\cap\mathcal F(x_\infty)\neq\varnothing\)，则存在 \(\omega\in\mathcal F(x_\infty)\) 使 \(\tau_{\infty\to m}(\omega)=a\)，并满足 \(\Fold_m(a)=x_m\)；
故所有可达前缀集合包含于纤维 \(\Fold_m^{-1}(x_m)\)，得到 \(N_m(x_\infty)\le d_m(x_m)\)。
\par
对复杂度上界：给定 \((x_\infty,m)\) 可恢复 \(x_m=\pi_m(x_\infty)\)。
由 \(\Fold_m\) 可计算且纤维有限可枚举，可对集合 \(\Fold_m^{-1}(x_m)\) 取字典序枚举并用秩 \(i\in\{0,\dots,d_m(x_m)-1\}\) 编码元素；
该秩可用 \(\lceil\log_2 d_m(x_m)\rceil\) 比特描述，而解释器与枚举程序只引入 \(O(\log m)\) 的开销。
故得所示上界；最后一式由两边除以 \(m\) 并取 \(\liminf\) 得到。
证毕。
\end{proof}

\begin{proposition}[路径分量纤维的 rank/unrank 线性时间实现]\label{prop:fold-fiber-rank-unrank-linear}
对整数 \(\ell\ge 0\)，令
\[
\mathcal{L}_\ell:=\{s\in\{0,1\}^{\ell}:\ s\ \text{不含相邻子串 }11\},
\qquad
\card{\mathcal{L}_\ell}=F_{\ell+2}.
\]
存在确定性算法 \(\mathrm{rank}_\ell:\mathcal{L}_\ell\to[0,F_{\ell+2}-1]\) 与 \(\mathrm{unrank}_\ell:[0,F_{\ell+2}-1]\to\mathcal{L}_\ell\)，互为逆，并且二者均在 \(O(\ell)\) 时间内可计算。
进一步，若某纤维在某一规范双射下分解为有限直积
\[
\Fold_m^{-1}(x)\ \cong\ \prod_{i=1}^{r}\mathcal{L}_{\ell_i},
\qquad
d_m(x)=\prod_{i=1}^{r}F_{\ell_i+2},
\]
则可用混合进制分解把纤维内的 rank/unrank 归约为分量 rank/unrank，从而在 \(O(\sum_i\ell_i)=O(m)\) 时间内实现纤维内的 rank/unrank。
\end{proposition}
\begin{proof}
对 \(\mathcal{L}_\ell\) 的 rank/unrank 用 Fibonacci 递推实现。
对任意 \(s=s_1\cdots s_\ell\in\mathcal{L}_\ell\)，若 \(s_1=0\)，则其在字典序下与 \(\mathcal{L}_{\ell-1}\) 的尾串秩一致；若 \(s_1=1\)，则必有 \(s_2=0\)，并且以 \(0\) 起始的合法串共有 \(\card{\mathcal{L}_{\ell-1}}=F_{\ell+1}\) 个，故
\[
\mathrm{rank}_\ell(s)=F_{\ell+1}+\mathrm{rank}_{\ell-2}(s_3\cdots s_\ell).
\]
反秩同理：给定 \(j\in[0,F_{\ell+2}-1]\)，若 \(j<F_{\ell+1}\) 则取首位为 \(0\) 并递归于 \(\ell-1\)；否则取前两位为 \(10\) 并递归于 \(\ell-2\) 的余数 \(j-F_{\ell+1}\)。
每步长度减少 \(1\) 或 \(2\)，故时间为 \(O(\ell)\)。
\par
对直积情形，令 \(N_i:=F_{\ell_i+2}\)。
任意 \(j\in[0,\prod_i N_i-1]\) 唯一分解为混合进制坐标
\(
j=\sum_i j_i\prod_{k<i}N_k
\)
（其中 \(0\le j_i<N_i\)）。
对每个分量执行 \(\mathrm{unrank}_{\ell_i}(j_i)\) 并拼接即得全局 unrank；rank 由逐分量 rank 得到 \((j_i)\) 后再作同一混合进制合成。
总时间为 \(O(\sum_i\ell_i)=O(m)\)。证毕。
\end{proof}

\begin{proposition}[终止性、局部合流与合流（纽曼引理）]\label{prop:fold-rewrite-newman}
上述重写系统在 $\mathcal{D}$ 上强终止，且局部合流；由纽曼引理可得其合流性。因此每个 $a\in\mathcal{D}$ 的正规形存在且唯一，记为 $\mathrm{NF}(a)\in X_\infty$，并满足
\[
\mathrm{Val}(\mathrm{NF}(a))=\mathrm{Val}(a).
\]
\leanverified{step\_stronglyTerminating}
\leanverified{step\_confluent}
\leanverified{paper\_fold\_rewrite\_newman}
\end{proposition}

\begin{proof}
\textbf{（终止性）}令 $A(a)=\sum_{k\ge 1} a_k$、$B(a)=\sum_{k\ge 1} k\,a_k$、$C(a)=a_2$，并设 $\mu(a)=(A(a),B(a),C(a))$ 按字典序比较。则规则 (\ref{rule:fold:adj}) 与 (\ref{rule:fold:dedup1}) 使 $A$ 减少；规则 (\ref{rule:fold:dedupk})（$k\ge 3$）使 $B$ 减少；规则 (\ref{rule:fold:dedup2}) 保持 $A,B$ 不变但使 $C$ 减少。故每步重写都严格降低 $\mu$，系统强终止。

\textbf{（局部合流）}若 $a$ 有两条一步归约 $a\to a_1,a_2$，则两侧继续归约到不可约形后都落在 $X_\infty$，且由于每步保持 $\mathrm{Val}$，两不可约形同值；由 Zeckendorf 唯一性它们相同，于是两支可汇合。

\textbf{（合流）}由纽曼引理（强终止 + 局部合流 $\Rightarrow$ 合流）即得结论。
\end{proof}

\begin{corollary}[\texorpdfstring{$\Fold_m$}{Fold\_m} 的归约顺序无关实现]\label{cor:foldm-order-indep}
对任意 $\omega\in\Omega_m$，
\[
\Fold_m(\omega)\ =\ \pi_m\bigl(\mathrm{NF}(\iota_m(\omega))\bigr),
\]
且右端与归约时选取规则/位置的顺序无关。
\leanverified{normalPrefix\_iota\_eq\_Fold}
\leanverified{paper\_foldm\_order\_indep}
\end{corollary}

\begin{remark}[规范截面与可计算性]\label{rem:fold-normal-form-computable-section}
命题 \ref{prop:fold-rewrite-newman} 给出一个可计算的规范截面：对任意 $a\in\mathcal{D}$，正规形 $\mathrm{NF}(a)\in X_\infty$ 唯一存在且与归约顺序无关。因此，当“第三参数”仅对应于等价类中代表的选择（规范冗余）时，无需显式记录该参数；可通过求正规形在多项式时间内选出唯一代表（分辨率口径为线性，见命题 \ref{prop:fold-complexity}）。需要强调的是，该消去仅恢复稳定代表，不恢复同一纤维内的微态残差。
\end{remark}

\begin{theorem}[折叠的构造--刚性--障碍三分]\label{thm:fold-constructibility-rigidity-obstruction}
\leanverified{paper\_fold\_constructibility\_rigidity\_obstruction}
固定分辨率 \(m\ge 1\) 与稳定词 \(x\in X_m\)，记
\[
d_m(x):=\bigl|\Fold_m^{-1}(x)\bigr|.
\]
则折叠 \(\Fold_m\) 的对象论地位可分为三层：
\begin{enumerate}
  \item \textbf{构造层：}每个稳定词 \(x\) 都有一个显式 admissible witness，事实上 \(x\in\Omega_m\) 且 \(\Fold_m(x)=x\)。
  \item \textbf{刚性层：}若 \(a,b\in\mathcal D\) 满足 \(\mathrm{Val}(a)=\mathrm{Val}(b)\)，则
  \[
  \pi_m(\mathrm{NF}(a))=\pi_m(\mathrm{NF}(b)).
  \]
  换言之，分辨率 \(m\) 的稳定代表只依赖于值类，而不依赖于局部重写路径或初始冗余表示。
  \item \textbf{障碍层：}从 \(x\) 恢复唯一微态当且仅当 \(d_m(x)=1\)。
  若 \(d_m(x)>1\)，则以 \(\pi=\delta_x\) 代入定理 \ref{thm:fold-microstate-residual-window-reachability}，可实现的纤维残差恰覆盖闭区间
  \[
  [0,\log d_m(x)],
  \]
  等价地，可实现的纤维熵恰覆盖同一区间。
\end{enumerate}
\end{theorem}
\begin{proof}
\textbf{(1)} 由命题 \ref{prop:fold-basic} 的幂等性，任意 \(x\in X_m\subset\Omega_m\) 都满足 \(\Fold_m(x)=x\)，故稳定对象先以显式输入的形式被合法构造出来。

\textbf{(2)} 由命题 \ref{prop:fold-rewrite-newman}，正规形 \(\mathrm{NF}(a)\in X_\infty\) 唯一存在并仅由 \(\mathrm{Val}(a)\) 决定；故若 \(\mathrm{Val}(a)=\mathrm{Val}(b)\)，则 \(\mathrm{NF}(a)=\mathrm{NF}(b)\)，再取前 \(m\) 位即得
\(
\pi_m(\mathrm{NF}(a))=\pi_m(\mathrm{NF}(b))
\)。
这与推论 \ref{cor:foldm-order-indep} 一致地说明：折叠读出在值类上是刚性的。

\textbf{(3)} 若 \(d_m(x)=1\)，则纤维 \(\Fold_m^{-1}(x)\) 仅含一个元素，故微态恢复唯一。
反之若 \(d_m(x)>1\)，则唯一恢复显然失败。
将定理 \ref{thm:fold-microstate-residual-window-reachability} 中的宏观分布取为 \(\pi=\delta_x\)，便得
\[
\{\mathsf R(\mu):\ (\Fold_m)_{\#}\mu=\delta_x\}=[0,\log d_m(x)],
\]
且对应的熵谱为
\[
\{H(\mu):\ (\Fold_m)_{\#}\mu=\delta_x\}=[0,\log d_m(x)].
\]
因此当纤维非单点时，障碍不是“没有稳定对象”，而是存在一个不可由 \(x\) 本身消去的纤维残差窗口。
证毕。
\end{proof}

\begin{theorem}[纤维残差窗口的全区间可达性与熵谱完备]\label{thm:fold-microstate-residual-window-reachability}
设 \(\Omega\) 为有限集合，\(\Fold:\Omega\twoheadrightarrow X\) 为满射。
对每个 \(x\in X\) 记纤维 \(F_x:=\Fold^{-1}(x)\) 与多重度 \(d(x):=|F_x|\)。
给定任意宏观分布 \(\pi\in\Delta(X)\)，定义其\textbf{纤维均匀提升} \(\mu^{\mathrm{e}}_\pi\in\Delta(\Omega)\) 为
\[
\mu^{\mathrm{e}}_\pi(\omega):=\frac{\pi(x)}{d(x)},\qquad \omega\in F_x.
\]
并令
\[
\mathcal{C}_\pi:=\{\mu\in\Delta(\Omega):\ \Fold_{\#}\mu=\pi\}.
\]
定义相对于纤维均匀提升的\textbf{微态残差}
\[
\mathsf{R}(\mu):=D_{\mathrm{KL}}\!\left(\mu\ \middle\|\ \mu^{\mathrm{e}}_\pi\right),
\]
以及\textbf{可逆窗口长度}
\[
W_{\mathrm{inv}}(\pi):=\sum_{x\in X}\pi(x)\log d(x).
\]
则成立全区间可达性：
\[
\boxed{\ \{\mathsf{R}(\mu):\ \mu\in\mathcal{C}_\pi\}=[0,W_{\mathrm{inv}}(\pi)]\ }.
\]
等价地，\(\mathcal{C}_\pi\) 中可实现的 Shannon 熵恰覆盖闭区间
\[
\boxed{\ \{H(\mu):\ \mu\in\mathcal{C}_\pi\}=\bigl[H(\pi),\ H(\pi)+W_{\mathrm{inv}}(\pi)\bigr]\ }.
\]
其中右端点由 \(\mu=\mu^{\mathrm{e}}_\pi\) 取到（纤维均匀提升为最大熵 lift），左端点由每条纤维取原子条件分布取到。
\end{theorem}
\begin{proof}
取任意 \(\mu\in\mathcal{C}_\pi\)，令 \(\mu_x\in\Delta(F_x)\) 表示 \(\mu\) 在纤维上的条件分布，则对 \(\omega\in F_x\) 有
\[
\mu(\omega)=\pi(x)\mu_x(\omega),
\qquad
\mu^{\mathrm{e}}_\pi(\omega)=\pi(x)u_x(\omega),
\]
其中 \(u_x\) 为 \(F_x\) 上的均匀分布。
由 KL 的链式分解（对有限分布为直接恒等式）得
\[
D_{\mathrm{KL}}\!\left(\mu\ \middle\|\ \mu^{\mathrm{e}}_\pi\right)
=\sum_{x\in X}\pi(x)\,D_{\mathrm{KL}}(\mu_x\|u_x).
\]
并且对每个 \(x\)，有 \(D_{\mathrm{KL}}(\mu_x\|u_x)=\log d(x)-H(\mu_x)\)，因而
\[
0\le D_{\mathrm{KL}}(\mu_x\|u_x)\le \log d(x),
\]
且当 \(\mu_x\) 在 \(\Delta(F_x)\) 上变化时，\(H(\mu_x)\) 连续覆盖区间 \([0,\log d(x)]\)，从而 \(D_{\mathrm{KL}}(\mu_x\|u_x)\) 连续覆盖 \([0,\log d(x)]\)。
因此
\[
\{\mathsf{R}(\mu):\mu\in\mathcal{C}_\pi\}
=\sum_{x\in X}\pi(x)\,[0,\log d(x)]
=[0,\sum_{x\in X}\pi(x)\log d(x)]
=[0,W_{\mathrm{inv}}(\pi)],
\]
其中“和”按 Minkowski 和理解；最后一步由区间加法闭包性直接给出。
\par
另一方面，直接计算得
\[
\mathsf{R}(\mu)=D_{\mathrm{KL}}\!\left(\mu\ \middle\|\ \mu^{\mathrm{e}}_\pi\right)
=H(\mu^{\mathrm{e}}_\pi)-H(\mu),
\qquad
H(\mu^{\mathrm{e}}_\pi)=H(\pi)+W_{\mathrm{inv}}(\pi),
\]
从而熵谱结论与残差谱结论互相等价。证毕。
\end{proof}

\endinput


\begin{proposition}[折叠的规范性、幂等性与满射性]\label{prop:fold-basic}
折叠算子 $\Fold_m$ 良定义且满足：
\begin{enumerate}
\normalfont
  \item \textbf{（规范性）} $\Fold_m(\omega)\in X_m$；
  \item \textbf{（幂等性）} 若 $\omega\in X_m$，则 $\Fold_m(\omega)=\omega$；
  \item \textbf{（满射性）} $\Fold_m:\Omega_m\twoheadrightarrow X_m$ 为满射。
\end{enumerate}
\leanverified{Fold\_stable}
\leanverified{Fold\_idempotent}
\leanverified{Fold\_surjective}
\leanverified{paper\_fold\_surjective}
\leanverified{stableValue\_Fold\_lt}
\leanverified{exactWeightCount\_one\_eq}
\leanverified{paper\_fold\_basic}
\leanverified{paper\_allTrue\_not\_no11}
\end{proposition}

\begin{proof}
\textbf{（规范性）}由定义 \ref{def:fold-word}，$Z(N(\omega))\in X_\infty$，故 $\pi_m(Z(N(\omega)))\in X_m$。

\textbf{（幂等性）}若 $\omega\in X_m$，令 $N=\sum_{k=1}^m \omega_k F_{k+1}$。由于 $\omega$ 已是合法 Zeckendorf 数字的长度-$m$ 前缀，且 Zeckendorf 表示唯一，故 $Z(N)$ 的前 $m$ 位必等于 $\omega$，从而 $\Fold_m(\omega)=\omega$。

\textbf{（满射性）}对任意 $x\in X_m$，取 $\omega=x\in\Omega_m$，由幂等性得 $\Fold_m(\omega)=x$，故 $\Fold_m$ 满射。
\end{proof}
